\section{Introduction}\label{sec:intro}

Visual language is the system that uses tightly integrated textual and visual elements to convey meaning \citep{horn1998visual}.
It is ubiquitous in the human world with typical examples being charts, plots and diagrams existing in places such as textbooks, scientific papers web pages and many more. % infographics etc.
Visual language is also highly complex -- besides texts, its structural units can include line, shape, color, orientation, scale, angle, space, etc. One needs to recognize patterns from these structural units, and perform spatial grouping and/or alignment to extract information for reasoning.

Whilst being prevalent and important, there is little research on visual language understanding from the machine learning community.
%has been largely overlooked by the machine learning community.
Vision-language models pretrained on natural images or image-text pairs crawled from the web perform badly on visual language tasks such as ChartQA \citep{masry-etal-2022-chartqa} and PlotQA \citep{methani2020plotqa} due to the high complexity of jointly modeling language and symbols (more evidence in experiments). Pix2Struct \citep{lee2022pix2struct} is a recently proposed pretraining strategy for visually-situated language that significantly outperforms standard vision-language models, and also a wide range of OCR-based pipeline approaches. Pix2Struct designs a novel masked webpage screenshot parsing task and also a variable-resolution input representation for pretraining an image-to-text encode-decoder Transformer \citep{vaswani2017attention}. In this work, we use Pix2Struct as the base model and further pretrain it with chart derendering and math reasoning tasks.

\begin{figure*}[!ht]
\centering
\includegraphics[width=\linewidth]{figs/matcha_front_page_fig.pdf}
\caption{\model{} defines two types of pretraining tasks: (1) chart derendering (light blue boxes) and (2) mathematical reasoning (light red boxes). In chart derendering, given a chart, the model needs to decode its underlying rendering code or data table. In math reasoning, given a math question rendered as an image, the model needs to decode its answer. Chart derendering teaches the models layout understanding (including number extraction and their organizations) and math reasoning teaches the models numerical reasoning capabilities. }
\label{fig:front_page}
\end{figure*}

We argue that visual language understanding needs two key ingredients: (1) layout understanding (including number extraction and their organizations) and (2) mathematical reasoning. (1) is required to discover the underlying patterns of the image and organize the elements in the image in a logical form. (2) is needed to operate on the elements extracted from (1) and derive meaningful information demanded by a task or query. Based on these observations, we propose two complementary pretraining tasks for enhancing visual language understanding: \textbf{chart derendering} and \textbf{math reasoning}. In chart derendering, given a plot/chart, the image-to-text model is required to generate its underlying data table or the code used to render it. The second task is math reasoning pretraining. We pick two numerical reasoning dataset MATH \citep{saxton2019analysing} and DROP \citep{dua-etal-2019-drop}, render the input into images and the image-to-text model needs to decode the answers.


We use a suite of visual language tasks to test the effectiveness of our method. Most importantly, we test on ChartQA and PlotQA which are QA datasets about plots and charts. On both datasets, \model{} surpasses even the SOTA model assuming access to charts' underlying data tables and can beat the prior SOTA without gold data table by as much as 20\%. We also test \model{} on chart-to-text summarization tasks and observe clear improvements over Pix2Struct and achieves SOTA on Chart-to-Text \citep{kantharaj-etal-2022-chart} Pew split. Last but not least, to examine if the \model{} pretraining generalizes to datasets beyond the standard plots and charts domain, we also test \model{} on four additional domains where Pix2Struct was evaluated on: documents, illustrations, user
interfaces, and natural images (including datasets, such as textbook QA, Widget Captioning, etc.).
We demonstrate consistent improvement on most additional datasets compared with the base model Pix2Struct. 

To summarize, our contributions are: (1) proposing a set of effective pretraining tasks for visual language learning (2) demonstrating consistent improvements across all evaluated tasks and SOTA results on ChartQA, PlotQA, and Chart-to-Text summarization (Statista set) without accessing the gold data tables; (3) verify that \model{} pretraining transfers to visual language benchmarks beyond the chart \& plot domains and achieve SOTA across a wide range of datasets beyond the chart domain such as textbook VQA and Widget Captioning; (4) comprehensive ablation and analyses to understand the effect of each pretraining component and its impact to downstream performance.
